\section{Reference-calibrated Attribution with Abstention}\label{app:selective}
This retrospective test asks when an existing reference-prototype assignment
should be issued rather than withheld. Its rule is fixed using historical
artworks; no generated prompt label fits or tunes the gate. The same generated
images and their earlier baseline confusions were previously inspected, so
the experiment is not independent confirmation. The primary condition is CSD
with original primary prototypes. Both encoders and every original/audited
view and primary/development prototype target remain mandatory outputs.

\subsection{Fixed calibration and decision rule}
Use the same unit-normalized historical prototype directions $u_a$ as the
existing baseline. The other historical partition supplies calibration:
development works for primary prototypes, primary works for development
prototypes. The selected source view is used consistently in both partitions;
an audited crop replaces only the previously declared region of that work.

For a calibration artwork $x$ labeled $a$, define nonconformity
\begin{equation}
 s_a(x)=\max_{b\ne a}x^\top u_b-x^\top u_a.
\end{equation}
With $n_a$ calibration works for painter $a$, set
$k_a=\lceil(n_a+1)(1-.10)\rceil$ and take $q_a$ to be the $k_a$-th smallest
score, using $+\infty$ when $k_a>n_a$. There is no interpolated quantile or
alternative nominal level. For each generated query form
\begin{equation}
 C(x)=\{a:s_a(x)\le q_a\},\qquad
 \hat a(x)=\arg\max_a x^\top u_a.
\end{equation}
Issue the unchanged top-1 assignment only if $C(x)=\{\hat a(x)\}$.
Otherwise abstain, retaining empty, multiple-label and disagreeing-singleton
sets as separate reasons. Ties in top-1 follow the fixed painter order.

This uses an ordinary rank-calibration rule as an empirical diagnostic.
Exchangeability between historical art and generated images is not established;
the numerical $.10$ choice is not a guarantee of generated-image coverage or
selective error. No density-ratio correction or shift assumption from
\citet{tibshirani2019shift} is fitted. The general rejection/coverage tradeoff
is established work \citep{geifman2017selective}, not an algorithmic novelty.

\subsection{Coverage-matched comparison and operational criterion}
Within a configuration's 112 named queries, let $K$ be the number accepted by
the reference gate. The comparator keeps exactly the $K$ greatest top-two
prototype margins, breaking ties by frozen observation ID. It uses no true
prompt name in selection. This is a transductive descriptive comparator: its
cutoff uses the query distribution, whereas the reference gate's thresholds
use historical artworks. Predictions remain the original top-1 assignments.
For controls, apply the resulting numeric margin cutoff inclusively, retaining
boundary ties. If $K=0$, neither named queries nor controls are accepted; if
$K=112$, the minimum named margin is still the cutoff for controls. No prompt
label determines either selection rule.

The pre-outcome operational criterion requires at least 50\% named-image
coverage in every configuration, at least one accepted query from every
prompted painter in each configuration, and positive equal-configuration mean
reductions in accepted error relative to both unrestricted top-1 and the
coverage-matched comparator. These joint descriptive conditions prevent a
nearly empty accepted set from being credited as useful. Undefined accepted
error is retained as a failure, without dropping its configuration from an
otherwise favorable mean. All individual adverse results remain visible.
The 50\% threshold is this study's declared operating requirement, not a
universal criterion for useful selective classification.

For every encoder/view/target/configuration, retain the individual scores,
candidate sets, predictions, decisions, painter-specific coverage, accepted
and abstained error, and correct/incorrect counts. Delete each scene once to
describe sensitivity, without changing historical calibration or interpreting
the ranges as sampling intervals. Control-arm results report name-assignment
rates separately for artist-free and generic-painting images. These rates
concern assignments without a candidate name in the prompt; they do not
establish that such an image has no resemblance to that painter.
